\documentclass[11pt,a4paper]{article}
\usepackage[T1]{fontenc}
\usepackage[utf8]{inputenc}
\usepackage{amsmath,amssymb,xcolor}
\usepackage[margin=25mm]{geometry}
\usepackage[colorlinks=true,urlcolor=blue]{hyperref}
\definecolor{nexorablue}{HTML}{2459B5}
\setlength{\parindent}{0pt}
\setlength{\parskip}{7pt}
\setlength{\emergencystretch}{2em}
\renewcommand{\familydefault}{\sfdefault}
\begin{document}

{\small\bfseries\color{nexorablue} NEXORA CPGE}\par
{\small MPSI / NEX-RGT-06 / Corrigé}\par
{\LARGE\bfseries Bilan chiffré d’une charge complète}\par
\vspace{0.25em}\hrule\vspace{0.6em}

\section*{Bilan numérique}
La source fournit $E(CE)=CE^2\simeq17{,}8\,\mathrm{mJ}$. Il reste $CE^2/2\simeq8{,}91\,\mathrm{mJ}$ dans la capacité. La différence, soit $8{,}91\,\mathrm{mJ}$, est dissipée.

\section*{Influence de la résistance}
Doubler $R$ double $RC$ et divise le courant initial par deux. Les états initial et final sont identiques et les trois énergies calculées sont les mêmes.


\end{document}
